\documentclass[11pt]{article}

% ============================================================
%  PLANTILLA ESTILO "IOC / EAC"
%  Edita el bloque "DATOS DEL DOCUMENTO" y los interruptores de
%  abajo, y luego el contenido a partir de \begin{document}.
% ============================================================

\usepackage[a4paper,left=2cm,right=2cm,top=2.2cm,bottom=3.5cm,headheight=28pt,footskip=30pt]{geometry}
\usepackage[utf8]{inputenc}
\usepackage{xcolor}
\usepackage{tabularx}
\usepackage{multirow}
\usepackage{array}
\usepackage{fancyhdr}
\usepackage{enumitem}
\usepackage{lastpage}
\usepackage{graphicx}

\definecolor{iocgray}{gray}{0.85}     % gris caja de título / barras de sección
\definecolor{iocbarragris}{gray}{0.80}
\definecolor{iocblau}{RGB}{0,102,204} % azul para "Nom i cognoms"

% ------------------------------------------------------------
%  INTERRUPTORES (activa/desactiva partes opcionales)
% ------------------------------------------------------------
\newif\ifiocfooter  \iocfootertrue   % \iocfooterfalse  -> sin pie de página
\newif\ifiocindex   \iocindexfalse   % \iocindextrue    -> añade índice tras la portada
\newif\ifiocbib     \iocbibfalse     % \iocbibtrue      -> añade bibliografía al final

% ------------------------------------------------------------
%  DATOS DEL DOCUMENTO  (cambia esto en cada nuevo EAC)
% ------------------------------------------------------------
\newcommand{\nomAlumne}{Nom i cognoms}
\newcommand{\ciclePrograma}{CEFP Ciberseguretat en Entorns de les Tecnologies de la Informació}
\newcommand{\modulNom}{Mòdul 5021 -- Incidents de ciberseguretat}
\newcommand{\codiActivitat}{EAC1}
\newcommand{\cursSemestre}{Curs 2026--27 / 1r semestre}
\newcommand{\subtitolActivitat}{Proposta d'awareness training}

\newcommand{\codiDocument}{I71}
\newcommand{\nomExercici}{Exercici d'avaluació contínua 1}
\newcommand{\versioDocument}{03}
\newcommand{\fitxerDocument}{CECIB\_5021\_EAC1\_Enunciat\_2627S1 -- Jordi Foixench}
\newcommand{\dataLliurament}{01/10/2026}

% ------------------------------------------------------------
%  PIE DE PÁGINA (tabla de metadatos; se define una vez y se
%  reutiliza, controlada por el interruptor \ifiocfooter)
% ------------------------------------------------------------
\newcommand{\peuEAC}{%
  \ifiocfooter
  \renewcommand{\arraystretch}{1.15}
  \noindent
  \begin{tabularx}{\textwidth}{|>{\centering\arraybackslash}p{1.6cm}|>{\centering\arraybackslash}p{1.4cm}|X|>{\centering\arraybackslash}p{2.6cm}|}
    \hline
    \multirow{2}{*}{\parbox[c]{1.5cm}{\centering \scriptsize\textbf{ioc}\\[-2pt]\tiny institut obert de catalunya}}
      & \scriptsize Codi: \codiDocument
      & \scriptsize \nomExercici
      & \scriptsize Pàgina \thepage\ de \pageref{LastPage} \\
    \cline{2-4}
      & \scriptsize Versió: \versioDocument
      & \scriptsize \fitxerDocument
      & \scriptsize Lliurament: \dataLliurament \\
    \hline
  \end{tabularx}
  \fi
}

% ------------------------------------------------------------
%  ESTILOS DE PÁGINA
%  - "eacportada": solo la portada (sin cabecera repetida, ya
%    que \portadaEAC dibuja su propia cabecera institucional).
%  - "eacpagina": el resto de páginas (cabecera con título +
%    nombre del alumno).
%  El pie de página es el mismo \peuEAC en ambos casos.
% ------------------------------------------------------------
\fancypagestyle{eacportada}{%
  \fancyhf{}
  \renewcommand{\headrulewidth}{0pt}
  \renewcommand{\footrulewidth}{0pt}
  \fancyfoot[C]{\peuEAC}
}

\fancypagestyle{eacpagina}{%
  \fancyhf{}
  \renewcommand{\headrulewidth}{0pt}
  \renewcommand{\footrulewidth}{0pt}
  \fancyhead[L]{\small \codiActivitat\ -- \subtitolActivitat}
  \fancyhead[R]{\small \textcolor{iocblau}{\nomAlumne}}
  \fancyfoot[C]{\peuEAC}
}

\pagestyle{eacpagina}

% ------------------------------------------------------------
%  COMANDOS DE ESTILO
% ------------------------------------------------------------

% Cabecera institucional + caja de título gris (solo en la portada)
\newcommand{\portadaEAC}{%
  \thispagestyle{eacportada}%
  \begin{minipage}[t]{0.7\textwidth}
    \small
    Generalitat de Catalunya\\
    Departament d'Educació\\
    \textbf{Institut Obert de Catalunya}
  \end{minipage}%
  \hfill
  \begin{minipage}[t]{0.28\textwidth}
    \raggedleft \textcolor{iocblau}{\nomAlumne}
  \end{minipage}

  \vspace{1.2cm}

  \noindent
  \colorbox{iocgray}{%
    \begin{minipage}{\dimexpr\textwidth-2\fboxsep\relax}
      \centering
      \vspace{0.4cm}
      \textbf{\ciclePrograma}\\[2pt]
      \textbf{\modulNom}\\[10pt]
      {\Huge \textbf{\codiActivitat}}\\[8pt]
      (\cursSemestre)
      \vspace{0.4cm}
    \end{minipage}%
  }
  \vspace{0.8cm}
}

% Barra gris de título de sección (p.ej. "Enunciat", "1. Introducció")
% También registra la entrada en el índice, si está activado.
\newcommand{\seccioEAC}[1]{%
  \vspace{0.4cm}
  \noindent\colorbox{iocbarragris}{%
    \begin{minipage}{\dimexpr\textwidth-2\fboxsep\relax}
      \textbf{#1}
    \end{minipage}%
  }
  \ifiocindex\addcontentsline{toc}{section}{#1}\fi
  \vspace{0.3cm}
}

% ============================================================
%  DOCUMENTO
% ============================================================
\begin{document}

\portadaEAC

% Índice opcional: activa \iocindextrue arriba para incluirlo.
\ifiocindex
  \renewcommand{\contentsname}{Índex}
  \tableofcontents
  \newpage
\fi

\seccioEAC{Enunciat}

\subsection*{\codiActivitat: \subtitolActivitat}

\textit{L'ús de la intel·ligència artificial en aquesta activitat pot afectar negativament
en l'aprenentatge i l'adquisició dels resultats d'aprenentatge. Cal prestar atenció a totes
les passes que aneu fent i mantenir l'esforç per tal que el vostre aprenentatge sigui òptim.}

En aquest EAC realitzareu un treball individual per simular una campanya de
conscienciació en ciberseguretat d'una empresa fictícia.

\seccioEAC{1. Introducció}

Text de la secció...

\seccioEAC{2. Activitat a realitzar}

Text de la secció...

\seccioEAC{3. Entrega}

\begin{enumerate}[label=\arabic*.]
  \item (9 punts) Captures de pantalla i explicacions del procés seguit.
  \item (6 punts) Explicació dels mètodes d'enganny utilitzats i proposta de mesures de
        prevenció de l'atac (des del punt de vista de la conscienciació dels empleats).
  \item (3 punts) Anàlisi dels resultats obtinguts.
  \item (3 punts) Conclusions sobre l'efectivitat de la campanya.
\end{enumerate}

% Bibliografía opcional: activa \iocbibtrue arriba para incluirla.
% Sustituye las entradas de ejemplo por las tuyas y cita con \cite{clave}.
\ifiocbib
  \renewcommand{\refname}{Bibliografia}
  \ifiocindex\addcontentsline{toc}{section}{Bibliografia}\fi
  \begin{thebibliography}{99}
    \bibitem{gophish} Gophish. \textit{Gophish: Open Source Phishing Framework}.
      \url{https://getgophish.com}

    \bibitem{enisa} ENISA. \textit{Cybersecurity awareness campaigns: practices and guidelines}.
      European Union Agency for Cybersecurity.
  \end{thebibliography}
\fi

\end{document}
